% ===== FILE 3/9: Tile rasterizer, Loss, PSNR, SSIM, Score tổng hợp =====

% --- Slide: Tile-based rasterizer ---
\begin{frame}{Kiến trúc rasterizer theo tile: vì sao render nhanh}
\scriptsize
Với hàng triệu Gaussian và hàng triệu pixel, việc lặp "cho mỗi pixel, quét qua mọi Gaussian" là bất khả thi về tốc độ.
3DGS giải quyết bằng kiến trúc \textbf{song song hoá theo tile}, chia màn hình thành các ô $16\times16$ pixel:
\begin{enumerate}\itemsep2pt
    \item \textbf{Xác định phạm vi:} với mỗi Gaussian, dùng compact box (slide trước) để biết nó phủ tới những tile
    nào $\Rightarrow$ sinh ra các cặp khoá $(\text{chỉ số tile}, \text{độ sâu})$ --- một Gaussian phủ nhiều tile thì sinh
    nhiều cặp khoá, mỗi cặp ứng với một tile.
    \item \textbf{Sắp xếp một lần duy nhất:} toàn bộ danh sách cặp khoá được sắp theo $(\text{tile}, \text{depth})$ ---
    đây là bước tốn kém nhất về mặt lý thuyết, nhưng chỉ thực hiện \textbf{một lần cho cả khung hình}, không phải một
    lần riêng cho từng pixel như cách làm ngây thơ.
    \item \textbf{Xử lý song song theo tile:} mỗi \emph{thread block} trên GPU phụ trách đúng một tile; vì danh sách đã
    sắp sẵn theo depth, mỗi thread chỉ cần blend tuần tự (công thức alpha blending, slide trước) theo đúng thứ tự, dừng
    sớm khi $T$ đủ nhỏ.
\end{enumerate}
\textbf{Kết quả:} tổng chi phí gần với $O(N\cdot K)$, với $K$ là số tile trung bình mà mỗi Gaussian chạm tới (bị chi
phối trực tiếp bởi \texttt{mult} của compact box) --- thay vì $O(N\cdot H\cdot W)$ nếu phải quét toàn ảnh cho mỗi
Gaussian. Đây là lý do cốt lõi giúp 3DGS render \textbf{thời gian thực}.
\end{frame}

% --- Slide: Hàm mất mát ---
\begin{frame}{Hàm mất mát: so khớp ảnh render với ảnh thật}
\scriptsize
Sau khi render ra ảnh $I_{\text{render}}$, cần một hàm đo "khác biệt" với ảnh thật $I_{gt}$ để lấy gradient huấn luyện.
Trong repo này (\texttt{train.py:256--259}):
\[
\boxed{\;\mathcal{L} = (1-\lambda_{\text{dssim}})\,L_1 + \lambda_{\text{dssim}}\,(1-\mathrm{SSIM}) + \lambda_{L2}\,L_2\;}
\]
với $\lambda_{\text{dssim}}=0.2$, $\lambda_{L2}=2.0$ --- so với 3DGS gốc (chỉ có hai số hạng đầu), bản này \textbf{cộng
thêm} một số hạng $L_2$ (bình phương sai số theo pixel, \texttt{utils/loss\_utils.py:25}).
\begin{itemize}\itemsep1pt
    \item $L_1=\overline{|I_{\text{render}}-I_{gt}|}$: ổn định, ít bị kéo lệch bởi vài pixel ngoại lệ (outlier) --- một
    vài pixel sai rất nhiều (do noise, occlusion...) không làm gradient tổng thể "hoảng loạn".
    \item $\mathrm{SSIM}$ (Structural Similarity, cửa sổ Gaussian $11\times11$, $\sigma=1.5$): so khớp \textbf{cấu trúc
    cục bộ} --- độ tương phản, độ chói theo từng vùng nhỏ --- gần với cảm nhận thị giác con người hơn so sánh từng pixel riêng lẻ.
    \item $L_2$ (bình phương): phạt \textbf{nặng hơn} với sai số lớn, giúp mô hình hội tụ nhanh ở giai đoạn đầu khi ảnh
    render còn rất khác ảnh thật.
    \item \textbf{Ghi chú quan trọng:} \texttt{loss\_utils.py} còn định nghĩa hai hàm \texttt{tone\_curve\_loss} và
    \texttt{frequency\_loss} (dựa trên structure tensor sẽ nói ở phần SADGS) --- nhưng \textbf{không có dòng nào trong
    \texttt{train.py} gọi tới hai hàm này}. Đây là \emph{dead code}: tồn tại trong file nhưng không ảnh hưởng gì tới
    loss thật sự đang được tối ưu.
\end{itemize}
\begin{center}
\includegraphics[width=0.34\linewidth]{ch5_loss_mix.png}
\end{center}
\end{frame}

% --- Slide: PSNR ---
\begin{frame}{Đánh giá chất lượng (1/3): PSNR}
\scriptsize
Sau khi train xong, cần các chỉ số \textbf{độc lập với loss} để đánh giá khách quan chất lượng render trên tập test
(ảnh chưa dùng để train). Chỉ số đầu tiên và đơn giản nhất là \textbf{PSNR} (Peak Signal-to-Noise Ratio):
\[
\mathrm{PSNR}=10\log_{10}\!\left(\frac{1}{\mathrm{MSE}}\right)\quad\text{(dB, với ảnh chuẩn hoá về }[0,1]\text{)}
\]
\begin{itemize}\itemsep1pt
    \item PSNR càng \textbf{cao} nghĩa là sai số bình phương trung bình theo pixel càng \textbf{thấp} --- là thang đo
    logarit nên chênh lệch vài dB đã tương ứng với chênh lệch sai số đáng kể.
    \item \textbf{Hạn chế:} PSNR đối xử với mọi pixel như nhau, không phân biệt "sai ở vùng chi tiết phức tạp" hay "sai ở
    vùng phẳng đơn giản" --- hai ảnh có cùng PSNR có thể trông rất khác nhau về mặt cảm nhận thị giác (một ảnh mờ đều,
    một ảnh có vài vệt lỗi rõ rệt nhưng phần còn lại rất sắc nét).
    \item Vì hạn chế này, PSNR thường được báo cáo \textbf{cùng với} SSIM và LPIPS (hai slide tiếp theo), không dùng một mình.
\end{itemize}
\begin{center}
\includegraphics[width=0.36\linewidth]{ch5_psnr.png}
\end{center}
\end{frame}

% --- Slide: SSIM ---
\begin{frame}{Đánh giá chất lượng (2/3): SSIM}
\scriptsize
\textbf{SSIM} (Structural Similarity Index) --- cùng công thức đã dùng làm một phần của loss (slide trước) --- nhưng ở
đây dùng để \textbf{đánh giá}, không phải để lấy gradient. So sánh ảnh theo \textbf{cửa sổ trượt cục bộ} (Gaussian
$11\times11$, $\sigma=1.5$), kết hợp ba thành phần: độ chói (luminance), độ tương phản (contrast), và cấu trúc
(structure) giữa hai ảnh trong mỗi cửa sổ.
\begin{itemize}\itemsep1pt
    \item Giá trị nằm trong $[-1,1]$ (thường báo cáo trong $[0,1]$ với ảnh tự nhiên), $1$ là khớp hoàn hảo.
    \item Vì so sánh theo \textbf{vùng cục bộ} thay vì từng pixel riêng lẻ, SSIM nhạy với sai lệch về \textbf{cấu trúc}
    (đường nét, hoạ tiết bị mờ/lệch) hơn là sai lệch nhỏ về độ sáng tổng thể --- gần với cách mắt người đánh giá "ảnh này
    có giống ảnh kia không" hơn PSNR thuần.
    \item Trong 3DGS, SSIM đóng vai trò \textbf{kép}: vừa là một phần của hàm loss lúc train (khuyến khích khớp cấu
    trúc), vừa là chỉ số đánh giá lúc test.
\end{itemize}
\begin{center}
\includegraphics[width=0.36\linewidth]{ch5_ssim_window.png}
\end{center}
\end{frame}

% --- Slide: Score tổng hợp ---
\begin{frame}{Đánh giá chất lượng (3/3): kết hợp nhiều chỉ số}
\scriptsize
Ba chỉ số PSNR, SSIM, và \textbf{LPIPS} (Learned Perceptual Image Patch Similarity --- khoảng cách giữa hai ảnh đo
trong không gian đặc trưng của một mạng CNN đã huấn luyện trước trên dữ liệu thị giác con người) đo những khía cạnh
\textbf{khác nhau và bổ sung cho nhau}:
\begin{itemize}\itemsep1pt
    \item PSNR: sai số pixel thô, dễ tính, dễ hiểu, nhưng ít tương quan với cảm nhận thị giác.
    \item SSIM: cấu trúc cục bộ, tương quan tốt hơn với cảm nhận, nhưng vẫn là một công thức toán học tường minh, không "học" từ dữ liệu thị giác người.
    \item LPIPS: học trực tiếp từ dữ liệu đánh giá của con người, bắt được các khác biệt "nhìn thấy rõ" mà PSNR/SSIM có
    thể bỏ sót (ví dụ: một chi tiết bị làm mờ nhẹ nhưng vẫn giữ đúng cấu trúc trung bình).
    \item Không chỉ số nào là "đủ một mình" --- một báo cáo kết quả nghiêm túc nên nêu \textbf{cả ba}, kèm rõ scene,
    checkpoint (số iteration), và cấu hình đo, để tránh so sánh khập khiễng giữa các thiết lập khác nhau.
\end{itemize}
\begin{center}
\includegraphics[width=0.36\linewidth]{ch5_score_decomp.png}
\end{center}
\end{frame}
